% ============================================================
% PHẦN 3: CƠ SỞ LÝ THUYẾT
% ============================================================
\section{Cơ sở lý thuyết}

% ------------------------------------------------------------
\subsection{Kiến trúc Kappa}

Kiến trúc Kappa được Jay Kreps đề xuất vào năm 2014 \cite{kreps2014kappa} như một
giải pháp thay thế đơn giản hơn cho kiến trúc Lambda. Trong khi Lambda Architecture
duy trì song song hai lớp xử lý --- Batch Layer (xử lý toàn bộ dữ liệu lịch sử
theo lô) và Speed Layer (xử lý dữ liệu mới nhất theo thời gian thực) --- kiến trúc
Kappa loại bỏ hoàn toàn Batch Layer và chỉ giữ lại một Speed Layer duy nhất dựa
trên hệ thống streaming.

\subsubsection{Nguyên lý hoạt động}

Kiến trúc Kappa xây dựng trên hai tiền đề cốt lõi. \textbf{Thứ nhất}, hệ thống
streaming hiện đại (điển hình là Apache Kafka) có khả năng lưu trữ toàn bộ dữ liệu
sự kiện lâu dài dưới dạng \emph{immutable log} --- nhật ký bất biến, có thể được
``replay'' (đọc lại từ đầu) bất cứ khi nào cần. \textbf{Thứ hai}, nếu một hệ thống
streaming có thể xử lý dữ liệu lịch sử cũng nhanh như dữ liệu mới phát sinh, thì
không cần thiết phải có một Batch Layer tách biệt.

Pipeline trong kiến trúc Kappa có cấu trúc tuyến tính:

\begin{figure}[H]
  \centering
  \setlength{\fboxsep}{12pt}
  \fbox{\begin{minipage}{0.88\textwidth}\centering\vspace{1.8cm}
    \FILL{Chèn hình tại đây}\\[0.6em]
    \textit{Gợi ý: Sơ đồ so sánh hai kiến trúc ---
    bên trái là Lambda Architecture (Batch Layer + Speed Layer + Serving Layer)
    và bên phải là Kappa Architecture (chỉ Speed Layer duy nhất). Dùng mũi tên
    và màu sắc để làm nổi bật sự đơn giản hóa của Kappa.
    Nguồn tham khảo: \url{https://www.oreilly.com/radar/questioning-the-lambda-architecture/}}
    \vspace{1.8cm}\end{minipage}}
  \caption{So sánh kiến trúc Lambda và kiến trúc Kappa}
  \label{fig:kappa-lambda}
\end{figure}

\begin{enumerate}[label=\arabic*.]
  \item \textbf{Nguồn dữ liệu (Data Sources):} Các thiết bị IoT, ứng dụng, cơ sở
    dữ liệu liên tục ghi sự kiện vào distributed log.
  \item \textbf{Distributed Log (Apache Kafka):} Đóng vai trò trung tâm lưu trữ
    và phân phối luồng sự kiện. Dữ liệu được giữ lại theo cấu hình
    \texttt{retention.ms}, cho phép consumer đọc lại lịch sử bằng cách reset
    offset về 0.
  \item \textbf{Speed Layer (Apache Spark Structured Streaming):} Tiêu thụ dữ liệu
    từ Kafka, áp dụng toàn bộ logic xử lý và tính toán kết quả, ghi kết quả vào
    Serving Layer.
  \item \textbf{Serving Layer (InfluxDB + Grafana):} Lưu trữ kết quả đã xử lý và
    phục vụ truy vấn từ người dùng hoặc dashboard.
\end{enumerate}

\subsubsection{Ưu điểm so với kiến trúc Lambda}

\begin{itemize}
  \item \textbf{Đơn giản hóa codebase:} Chỉ cần duy trì một luồng xử lý duy nhất,
    không còn tình trạng logic trùng lặp giữa Batch Layer và Speed Layer.
  \item \textbf{Dễ bảo trì:} Khi cần cập nhật thuật toán hoặc sửa lỗi, chỉ cần
    cập nhật một chỗ rồi replay lại dữ liệu từ Kafka.
  \item \textbf{Phù hợp cho dữ liệu IoT:} Dữ liệu cảm biến phát sinh theo thời
    gian thực với tần suất đều đặn, không cần xử lý lại theo lô hàng ngày.
\end{itemize}

Đề tài áp dụng kiến trúc Kappa cho pipeline BATADAL: Kafka đóng vai trò distributed
log lưu trữ dữ liệu cảm biến từ 43 điểm đo, Spark Structured Streaming thực hiện
toàn bộ tiền xử lý và phát hiện bất thường, InfluxDB lưu kết quả và Grafana hiển
thị dashboard thời gian thực.

\begin{figure}[H]
  \centering
  \setlength{\fboxsep}{12pt}
  \fbox{\begin{minipage}{0.88\textwidth}\centering\vspace{1.8cm}
    \FILL{Chèn hình tại đây}\\[0.6em]
    \textit{Gợi ý: Sơ đồ kiến trúc Kappa áp dụng cụ thể trong đề tài.
    Luồng từ trái sang phải: IoT Simulator (Python Producer) $\rightarrow$
    Apache Kafka (topic water\_stream) $\rightarrow$
    Spark Structured Streaming (rule-based + ML inference) $\rightarrow$
    InfluxDB (water\_telemetry) $\rightarrow$ Grafana Dashboard.
    Thêm mũi tên ``replay'' vòng từ Kafka về Spark để thể hiện
    tính năng xử lý lại lịch sử của kiến trúc Kappa.}
    \vspace{1.8cm}\end{minipage}}
  \caption{Kiến trúc Kappa áp dụng trong pipeline BATADAL của đề tài}
  \label{fig:kappa-pipeline}
\end{figure}

% ------------------------------------------------------------
\subsection{Apache Kafka}

Apache Kafka là một \emph{event streaming platform} phân tán, mã nguồn mở, được
phát triển bởi LinkedIn và chuyển giao cho Apache Software Foundation
\cite{kreps2011kafka}. Theo tài liệu chính thức tại \url{https://kafka.apache.org/intro},
Kafka kết hợp ba khả năng cốt lõi: (1) \emph{publish} và \emph{subscribe} luồng
sự kiện; (2) lưu trữ luồng sự kiện bền vững và đáng tin cậy trong thời gian dài;
và (3) xử lý luồng sự kiện theo thời gian thực hoặc hồi tố.

\subsubsection{Kiến trúc Kafka}

Kafka hoạt động như một hệ thống phân tán gồm \textbf{Servers} và \textbf{Clients}:

\begin{itemize}
  \item \textbf{Broker:} Kafka được vận hành như một cluster gồm một hoặc nhiều
    server (broker). Các broker chịu trách nhiệm nhận, lưu trữ và phục vụ dữ liệu
    sự kiện. Khi một broker gặp sự cố, các broker còn lại tự động tiếp quản để
    đảm bảo hoạt động liên tục không mất dữ liệu.
  \item \textbf{Producer:} Ứng dụng client ghi (publish) sự kiện vào Kafka topic.
    Producer quyết định partition nào sẽ nhận sự kiện, thường dựa trên event key.
  \item \textbf{Consumer:} Ứng dụng client đọc (subscribe) và xử lý sự kiện từ
    Kafka topic. Consumer theo dõi vị trí đọc bằng \emph{offset}, cho phép đọc
    lại từ bất kỳ điểm nào trong lịch sử.
\end{itemize}

\subsubsection{Khái niệm cốt lõi}

\begin{figure}[H]
  \centering
  \setlength{\fboxsep}{12pt}
  \fbox{\begin{minipage}{0.88\textwidth}\centering\vspace{1.8cm}
    \FILL{Chèn hình tại đây}\\[0.6em]
    \textit{Gợi ý: Sơ đồ kiến trúc tổng thể Apache Kafka thể hiện:
    (1) nhiều Producer ghi vào các Topic khác nhau;
    (2) Kafka Cluster gồm 3 Broker với các Partition và Replica;
    (3) nhiều Consumer Group độc lập đọc từ cùng một Topic;
    (4) ZooKeeper/KRaft quản lý metadata cluster.
    Hình tham khảo: \url{https://kafka.apache.org/intro} (phần ``How does Kafka work'').}
    \vspace{1.8cm}\end{minipage}}
  \caption{Kiến trúc tổng thể Apache Kafka: Broker, Producer và Consumer}
  \label{fig:kafka-arch}
\end{figure}

\textbf{Event} (sự kiện) là đơn vị dữ liệu cơ bản trong Kafka, gồm: event key,
event value, timestamp, và metadata headers tùy chọn.

\textbf{Topic} là nơi lưu trữ và tổ chức các sự kiện, tương tự như một thư mục
trong hệ thống file. Topic trong Kafka hỗ trợ nhiều producer ghi đồng thời và
nhiều consumer đọc độc lập. Khác với hệ thống message queue truyền thống, sự kiện
trong Kafka \emph{không bị xóa sau khi đọc} --- thời gian lưu trữ được cấu hình
qua \texttt{retention.ms}.

\textbf{Partition} là cơ chế phân mảnh dữ liệu: mỗi topic được chia thành nhiều
partition phân bố trên các broker khác nhau. Partition đảm bảo thứ tự sự kiện trong
phạm vi từng partition và cho phép đọc/ghi song song từ nhiều client. Sự kiện có
cùng key được ghi vào cùng một partition, đảm bảo thứ tự xử lý theo key.

\textbf{Replication factor} xác định số lượng bản sao dữ liệu duy trì trên các
broker khác nhau. Thiết lập phổ biến trong môi trường production là replication
factor = 3, đảm bảo tính sẵn sẵn sàng cao (high availability).

\begin{figure}[H]
  \centering
  \setlength{\fboxsep}{12pt}
  \fbox{\begin{minipage}{0.88\textwidth}\centering\vspace{1.8cm}
    \FILL{Chèn hình tại đây}\\[0.6em]
    \textit{Gợi ý: Sơ đồ chi tiết cấu trúc Topic -- Partition -- Offset trong Kafka.
    Vẽ một topic có 4 partition (P0, P1, P2, P3), mỗi partition là một
    danh sách sự kiện được đánh offset (0, 1, 2, \ldots).
    Hiển thị hai consumer thuộc hai consumer group khác nhau cùng đọc
    topic này với offset độc lập.
    Hình tham khảo: \url{https://kafka.apache.org/intro} (hình minh họa streams-and-tables-p1\_p4).}
    \vspace{1.8cm}\end{minipage}}
  \caption{Cấu trúc Topic, Partition và Offset trong Apache Kafka}
  \label{fig:kafka-topic}
\end{figure}

\subsubsection{KRaft Mode}

Từ phiên bản Kafka 3.3 trở lên, Kafka hỗ trợ chế độ KRaft (Kafka Raft) --- loại
bỏ hoàn toàn sự phụ thuộc vào Apache ZooKeeper để quản lý metadata và bầu leader.
Trong chế độ KRaft, Kafka tự quản lý cluster metadata thông qua giao thức đồng thuận
Raft nội bộ, giảm độ phức tạp vận hành và cải thiện thời gian khởi động cluster.
Đề tài này sử dụng Kafka phiên bản 8.1.2 (Confluent Platform) với KRaft mode, không
cần triển khai ZooKeeper riêng biệt.

% ------------------------------------------------------------
\subsection{Apache Spark Structured Streaming}

Apache Spark là một \emph{unified analytics engine} cho Big Data, hỗ trợ xử lý
cả batch và streaming trên cùng một nền tảng \cite{zaharia2016spark}. Spark
Structured Streaming, được giới thiệu trong Spark 2.0, là mô hình xử lý luồng dữ
liệu có khả năng chịu lỗi, xây dựng trên Spark SQL engine.

\subsubsection{Mô hình lập trình}

Theo tài liệu chính thức của Apache Spark
(\url{https://spark.apache.org/docs/latest/streaming/}), Structured Streaming cung
cấp mô hình lập trình \emph{``incremental query on unbounded table''}: luồng dữ liệu
vào được mô hình hóa như một bảng vô hạn (unbounded table) liên tục được bổ sung
thêm dòng mới. Người dùng viết truy vấn trên bảng này giống hệt như viết truy vấn
batch trên DataFrame/Dataset tĩnh --- Spark engine tự động chuyển đổi thành xử lý
streaming gia tăng (incremental processing).

\begin{figure}[H]
  \centering
  \setlength{\fboxsep}{12pt}
  \fbox{\begin{minipage}{0.88\textwidth}\centering\vspace{1.8cm}
    \FILL{Chèn hình tại đây}\\[0.6em]
    \textit{Gợi ý: Sơ đồ mô hình ``Unbounded Table'' của Spark Structured Streaming.
    Bên trái: luồng dữ liệu đến (input data stream) theo trục thời gian.
    Ở giữa: bảng vô hạn (unbounded table) liên tục thêm hàng mới.
    Bên phải: result table (bảng kết quả) được cập nhật sau mỗi trigger.
    Ở dưới: output sink (InfluxDB). Hình tham khảo:
    \url{https://spark.apache.org/docs/latest/structured-streaming-programming-guide.html}
    (phần ``Programming Model'').}
    \vspace{1.8cm}\end{minipage}}
  \caption{Mô hình Unbounded Table trong Spark Structured Streaming}
  \label{fig:spark-unbounded}
\end{figure}

\subsubsection{Micro-Batch Execution}

Chế độ mặc định của Structured Streaming là \emph{micro-batch execution}: engine
thu thập dữ liệu mới trong một khoảng thời gian (trigger interval), sau đó xử lý
toàn bộ micro-batch đó như một job Spark thông thường. Trigger interval có thể được
cấu hình linh hoạt:

\begin{itemize}
  \item \texttt{ProcessingTime}: xử lý theo chu kỳ cố định (ví dụ: mỗi 5 giây).
  \item \texttt{Once}: xử lý tất cả dữ liệu hiện có rồi dừng (phù hợp cho
    scheduled job).
  \item \texttt{Continuous}: chế độ thực nghiệm với độ trễ thấp ở mức
    mili-giây (row-at-a-time processing).
\end{itemize}

\subsubsection{Tích hợp Kafka}

\begin{figure}[H]
  \centering
  \setlength{\fboxsep}{12pt}
  \fbox{\begin{minipage}{0.88\textwidth}\centering\vspace{1.8cm}
    \FILL{Chèn hình tại đây}\\[0.6em]
    \textit{Gợi ý: Sơ đồ timeline minh họa cơ chế Micro-Batch Execution.
    Trục ngang là thời gian; đánh dấu các trigger interval đều nhau (ví dụ: 5s).
    Mỗi khoảng trigger là một micro-batch (tô màu khác nhau).
    Bên dưới mỗi micro-batch hiển thị: số record nhận được, thời gian xử lý.
    So sánh với mô hình Continuous Processing (row-at-a-time) phía dưới
    để làm nổi bật sự khác biệt.}
    \vspace{1.8cm}\end{minipage}}
  \caption{Cơ chế Micro-Batch Execution trong Spark Structured Streaming}
  \label{fig:spark-microbatch}
\end{figure}

Spark Structured Streaming tích hợp Kafka qua thư viện
\texttt{spark-sql-kafka-0-10}. Mỗi message Kafka được ánh xạ thành một row trong
DataFrame với các cột: \texttt{key}, \texttt{value}, \texttt{topic},
\texttt{partition}, \texttt{offset}, \texttt{timestamp}. Trường \texttt{value}
chứa payload dạng bytes, cần được deserialize (thường là JSON) trước khi xử lý.

\subsubsection{Pandas UDF và Broadcast Variable}

Để tích hợp mô hình scikit-learn vào Spark, đề tài sử dụng hai cơ chế:
\textbf{Broadcast variable} truyền mô hình đã huấn luyện đến mọi executor một lần
duy nhất, tránh gửi lại mô hình cho mỗi task. \textbf{Pandas UDF}
(\texttt{pandas\_udf}) cho phép áp dụng hàm Python thuần nhận và trả về
\texttt{pandas.Series} trực tiếp trên từng cột của DataFrame, khai thác Apache
Arrow để truyền dữ liệu hiệu quả giữa JVM và Python process.

% ------------------------------------------------------------
\subsection{InfluxDB}

InfluxDB là một \emph{time-series database} (TSDB) mã nguồn mở, được phát triển
bởi InfluxData và tối ưu cho việc ghi, lưu trữ và truy vấn dữ liệu có kèm timestamp
với tần suất cao. Đây là lựa chọn phổ biến cho các ứng dụng IoT, monitoring hệ
thống, và observability. Phiên bản InfluxDB 2.x được sử dụng trong đề tài.

\subsubsection{Mô hình dữ liệu}

Theo tài liệu chính thức tại
\url{https://docs.influxdata.com/influxdb/v2/reference/key-concepts/data-elements/},
InfluxDB 2.x tổ chức dữ liệu theo hệ thống phân cấp:

\begin{itemize}
  \item \textbf{Organization:} Đơn vị quản lý cấp cao nhất, nhóm các tài nguyên
    liên quan (bucket, dashboard, task, user).
  \item \textbf{Bucket:} Nơi lưu trữ dữ liệu time-series, tương đương với
    database trong các hệ RDBMS. Mỗi bucket có \emph{retention policy} quy định
    thời gian lưu dữ liệu trước khi tự động xóa.
  \item \textbf{Measurement:} Nhóm logic cho dữ liệu time-series, tương đương với
    table. Trong đề tài, measurement \texttt{water\_telemetry} chứa toàn bộ dữ liệu
    cảm biến.
  \item \textbf{Tags:} Cặp key-value lưu metadata, có giá trị thay đổi không
    thường xuyên (ví dụ: \texttt{sensor\_id}, \texttt{location}). Tags được lập
    chỉ mục (indexed) nên phù hợp cho lọc và nhóm dữ liệu.
  \item \textbf{Fields:} Cặp key-value lưu giá trị đo lường thực tế, thay đổi
    theo thời gian (ví dụ: mức nước, lưu lượng, áp suất). Fields không được lập
    chỉ mục.
  \item \textbf{Timestamp:} Mốc thời gian liên kết với mỗi điểm dữ liệu, độ chính
    xác nano-giây.
\end{itemize}

Mỗi \textbf{Point} (điểm dữ liệu) được xác định duy nhất bởi tổ hợp:
measurement + tag set + field key + timestamp. Tập hợp các điểm có cùng
measurement, tag keys và tag values tạo thành một \textbf{Series}.

\begin{figure}[H]
  \centering
  \setlength{\fboxsep}{12pt}
  \fbox{\begin{minipage}{0.88\textwidth}\centering\vspace{1.8cm}
    \FILL{Chèn hình tại đây}\\[0.6em]
    \textit{Gợi ý: Sơ đồ cây phân cấp mô hình dữ liệu InfluxDB 2.x.
    Từ trên xuống: Organization $\rightarrow$ Bucket (water\_bucket, retention) $\rightarrow$
    Measurement (water\_telemetry) $\rightarrow$
    Tags (sensor\_id, location) và Fields (L\_T1, F\_PU1, \ldots) và Timestamp.
    Ở dưới cùng vẽ một bảng ví dụ vài dòng dữ liệu thực từ water\_telemetry
    để minh họa cụ thể. Nguồn tham khảo:
    \url{https://docs.influxdata.com/influxdb/v2/reference/key-concepts/data-elements/}.}
    \vspace{1.8cm}\end{minipage}}
  \caption{Mô hình dữ liệu phân cấp của InfluxDB 2.x}
  \label{fig:influxdb-model}
\end{figure}

\subsubsection{Ngôn ngữ truy vấn Flux}

Flux là ngôn ngữ truy vấn và scripting chức năng (functional) được InfluxData phát
triển riêng cho InfluxDB 2.x. Flux sử dụng cú pháp \emph{pipeline} với toán tử
\texttt{|>} (pipe-forward), cho phép xây dựng chuỗi biến đổi dữ liệu một cách rõ
ràng:

\begin{lstlisting}[language=Python, caption={Ví dụ truy vấn Flux lấy dữ liệu mức nước trong 1 giờ qua}]
from(bucket: "water_bucket")
  |> range(start: -1h)
  |> filter(fn: (r) => r._measurement == "water_telemetry")
  |> filter(fn: (r) => r._field == "L_T1")
  |> aggregateWindow(every: 5m, fn: mean, createEmpty: false)
\end{lstlisting}

\subsubsection{Giao thức ghi dữ liệu}

InfluxDB sử dụng \emph{Line Protocol} --- định dạng văn bản thuần để ghi dữ liệu
qua HTTP API tại endpoint \texttt{/api/v2/write}. Mỗi dòng có cú pháp:

\begin{lstlisting}[caption={Cú pháp Line Protocol của InfluxDB}]
<measurement>[,<tag_key>=<tag_value>] <field_key>=<field_value> [<timestamp>]
\end{lstlisting}

Trong đề tài, thư viện \texttt{influxdb-client-python} được sử dụng để ghi dữ liệu
từ Spark Streaming vào InfluxDB thông qua \texttt{WriteApi} với chế độ
\texttt{SYNCHRONOUS}.

% ------------------------------------------------------------
\subsection{Grafana}

Grafana là một nền tảng trực quan hóa và quan sát (observability platform) mã nguồn
mở, phát triển bởi Grafana Labs. Theo tài liệu chính thức tại
\url{https://grafana.com/docs/grafana/latest/introduction/}, Grafana cho phép truy
vấn, trực quan hóa, cảnh báo và khám phá dữ liệu từ nhiều nguồn khác nhau thông
qua giao diện web thống nhất.

\subsubsection{Thành phần chính}

\textbf{Data Source:} Grafana hỗ trợ hàng chục loại nguồn dữ liệu, bao gồm
InfluxDB, Prometheus, Elasticsearch, PostgreSQL, MySQL và nhiều loại khác. Mỗi
data source được cấu hình với URL, thông tin xác thực và các tùy chọn kết nối
riêng.

\textbf{Panel:} Đơn vị trực quan hóa cơ bản trong Grafana. Mỗi panel thực hiện
truy vấn đến data source và hiển thị kết quả dưới nhiều dạng: time-series graph,
gauge, stat, bar chart, table, heatmap, và các dạng khác.

\textbf{Dashboard:} Tập hợp các panel được sắp xếp trên một canvas. Dashboard hỗ
trợ \emph{template variables} cho phép lọc dữ liệu động mà không cần chỉnh sửa
từng truy vấn.

\textbf{Alerting:} Grafana tích hợp cơ chế cảnh báo cho phép định nghĩa điều kiện
kích hoạt (alert rule) dựa trên giá trị dữ liệu. Khi điều kiện được thỏa mãn,
Grafana gửi thông báo qua email, Slack, PagerDuty hoặc webhook.

\begin{figure}[H]
  \centering
  \setlength{\fboxsep}{12pt}
  \fbox{\begin{minipage}{0.88\textwidth}\centering\vspace{1.8cm}
    \FILL{Chèn hình tại đây}\\[0.6em]
    \textit{Gợi ý: Chụp màn hình (screenshot) Grafana dashboard ``BATADAL Water Network''
    đang chạy thực tế (truy cập \texttt{http://localhost:3000}).
    Dashboard nên hiển thị: panel mức nước các bể $L_{T1}\ldots L_{T7}$ dạng
    time-series graph; panel trạng thái bơm dạng stat/gauge;
    panel cảnh báo rule-based và ML alert dạng table hoặc bar.
    Đây là minh chứng trực tiếp cho khả năng trực quan hóa thời gian thực
    của Grafana kết hợp InfluxDB.}
    \vspace{1.8cm}\end{minipage}}
  \caption{Giao diện Grafana dashboard giám sát mạng nước BATADAL theo thời gian thực}
  \label{fig:grafana-dashboard}
\end{figure}

\subsubsection{Grafana Provisioning}

Theo tài liệu chính thức tại
\url{https://grafana.com/docs/grafana/latest/administration/provisioning/},
Grafana Provisioning là cơ chế cấu hình tự động qua file YAML, cho phép:

\begin{itemize}
  \item \textbf{Datasource Provisioning:} Định nghĩa data source trong file
    \texttt{datasources/*.yaml}, Grafana tự động tạo và cấu hình khi khởi động.
  \item \textbf{Dashboard Provisioning:} Định nghĩa thư mục chứa file JSON
    dashboard trong \texttt{dashboards/*.yaml}, Grafana tự động import tất cả
    dashboard JSON tìm thấy.
\end{itemize}

Cơ chế này cho phép \emph{version control} toàn bộ cấu hình Grafana cùng với
source code, đảm bảo tính tái lập (reproducibility) khi triển khai hệ thống. Trong
đề tài, cấu hình datasource (InfluxDB) và dashboard (BATADAL water network) được
provisioning tự động thông qua Docker Compose volume mount.

% ------------------------------------------------------------
\subsection{Bộ dữ liệu BATADAL}

BATADAL (Battle of the Attack Detection ALgorithms) là bộ dữ liệu benchmark chuẩn
được xây dựng bởi Taormina và cộng sự \cite{taormina2018batadal} nhằm đánh giá và
so sánh các thuật toán phát hiện tấn công mạng vật lý (cyber-physical attack) trên
hệ thống mạng cấp nước (WDN).

\subsubsection{Mạng C-TOWN}

Dữ liệu BATADAL được tạo ra từ mô phỏng hệ thống C-TOWN --- một mạng nước đô thị
tổng hợp có quy mô vừa, được mô tả chi tiết trong file \texttt{CTOWN.INP} (định dạng
EPANET). C-TOWN gồm các thành phần:

\begin{itemize}
  \item 7 khu vực phân phối nước (district metered areas)
  \item 10 bể chứa nước ($T_1$ đến $T_{10}$, trừ $T_6$, $T_9$, $T_{10}$)
  \item 11 trạm bơm ($PU_1$ đến $PU_{11}$)
  \item 2 van điều tiết ($V_1$, $V_2$)
  \item Hàng nghìn nút tiêu thụ (demand nodes)
\end{itemize}

\subsubsection{Cấu trúc dữ liệu}

\begin{figure}[H]
  \centering
  \setlength{\fboxsep}{12pt}
  \fbox{\begin{minipage}{0.88\textwidth}\centering\vspace{1.8cm}
    \FILL{Chèn hình tại đây}\\[0.6em]
    \textit{Gợi ý: Sơ đồ topology mạng nước C-TOWN được vẽ từ file
    \texttt{CTOWN.INP} (dùng EPANET hoặc WNTR để xuất hình).
    Thể hiện: 7 khu vực phân phối (tô màu khác nhau), vị trí các bể chứa
    $T_1$--$T_7$ (hình trụ), các trạm bơm $PU_1$--$PU_{11}$ (hình tam giác),
    các van $V_1$, $V_2$ (hình thoi), và đường ống nối giữa các nút.
    Chú thích tên từng thành phần trực tiếp trên sơ đồ.}
    \vspace{1.8cm}\end{minipage}}
  \caption{Topology mạng nước C-TOWN --- nền tảng dữ liệu BATADAL}
  \label{fig:ctown-network}
\end{figure}

Bộ dữ liệu BATADAL ghi lại trạng thái hệ thống C-TOWN theo chu kỳ 1 giờ, với
43 đặc trưng (features) gồm ba nhóm:

\begin{itemize}
  \item $L_{Ti}$: mức nước bể chứa (tank level), đơn vị mét --- 7 giá trị.
  \item $F_{PUj}$ và $S_{PUj}$: lưu lượng (flow) và trạng thái bật/tắt (status)
    của từng máy bơm --- mỗi loại 11 giá trị.
  \item $F_{V1}$, $S_{V2}$: lưu lượng qua van $V_1$ và trạng thái van $V_2$.
  \item Timestamp theo định dạng \texttt{DD/MM/YY HH}.
  \item \texttt{ATT\_FLAG}: nhãn tấn công ($0$ = bình thường, $1$ = tấn công,
    $-999$ = chưa gán nhãn).
\end{itemize}

\subsubsection{Các tập dữ liệu}

\begin{itemize}
  \item \textbf{Dataset 1} (training set): 8\,761 bản ghi từ 01/01/2016 đến
    31/12/2016 (\texttt{ATT\_FLAG} = 0 toàn bộ --- không có tấn công). Dùng
    để học mô hình trạng thái bình thường.
  \item \textbf{Dataset 2} (test set with labels): 4\,393 bản ghi với nhãn
    tấn công đầy đủ (\texttt{ATT\_FLAG} $\in \{0,1\}$). Dùng để đánh giá mô hình
    phát hiện bất thường.
  \item \textbf{Dataset 3} (test set without labels): tương tự Dataset 2 nhưng
    \texttt{ATT\_FLAG} = $-999$ (dữ liệu mù). Trong đề tài, Dataset 3
    (\texttt{BATADAL\_dataset04.csv}) được sử dụng để giả lập luồng dữ liệu
    streaming real-time.
\end{itemize}

\begin{figure}[H]
  \centering
  \setlength{\fboxsep}{12pt}
  \fbox{\begin{minipage}{0.88\textwidth}\centering\vspace{1.8cm}
    \FILL{Chèn hình tại đây}\\[0.6em]
    \textit{Gợi ý: Biểu đồ chuỗi thời gian (time-series plot) thể hiện
    giá trị mức nước bể $L_{T1}$ trong Dataset 2 trong một khoảng thời gian
    bao gồm cả giai đoạn bình thường và giai đoạn tấn công.
    Tô màu đỏ/hatch các khoảng thời gian có \texttt{ATT\_FLAG}=1.
    Trục X: timestamp; Trục Y: mức nước (mét).
    Hình này trực tiếp minh họa bản chất bài toán phát hiện bất thường
    trên chuỗi thời gian.}
    \vspace{1.8cm}\end{minipage}}
  \caption{Chuỗi thời gian mức nước $L_{T1}$: so sánh giai đoạn bình thường và tấn công}
  \label{fig:batadal-timeseries}
\end{figure}

Ngoài ra, đề tài xây dựng thêm \textbf{BATADAL\_synthetic.csv} --- tập dữ liệu
tấn công tổng hợp --- bằng cách áp dụng các kịch bản tấn công giả lập
(nhiễu Gaussian, can thiệp vào trạng thái bơm và mức nước bể) lên Dataset 1, nhằm
tạo ra lớp dương (positive class) cho quá trình huấn luyện mô hình học máy.

% ------------------------------------------------------------
\subsection{Random Forest}

Random Forest là thuật toán học máy thuộc nhóm phương pháp ensemble, được Leo
Breiman đề xuất năm 2001 \cite{breiman2001random}. Theo tài liệu chính thức của
scikit-learn tại
\url{https://scikit-learn.org/stable/modules/ensemble.html\#random-forests},
Random Forest xây dựng một tập hợp (ensemble) gồm nhiều cây quyết định (decision
tree) độc lập và tổng hợp kết quả dự đoán của chúng.

\subsubsection{Cơ chế hoạt động}

Random Forest kết hợp hai nguồn ngẫu nhiên để giảm phương sai (variance) của mô
hình đơn lẻ:

\begin{enumerate}[label=\arabic*.]
  \item \textbf{Bootstrap Sampling (Bagging):} Mỗi cây được huấn luyện trên một
    tập con dữ liệu được lấy mẫu ngẫu nhiên có hoàn lại (bootstrap sample) từ
    tập huấn luyện gốc. Xấp xỉ $\frac{1}{3}$ số mẫu không xuất hiện trong
    bootstrap sample của một cây gọi là \emph{out-of-bag (OOB) samples}, có thể
    dùng để ước lượng lỗi tổng quát hóa.
  \item \textbf{Random Feature Selection:} Tại mỗi nút chia (split node) trong
    quá trình xây dựng cây, chỉ một tập con ngẫu nhiên gồm $m$ đặc trưng
    (controlled by \texttt{max\_features}) được xem xét để chọn điểm chia tốt nhất.
    Mặc định scikit-learn dùng $m = \sqrt{p}$ cho bài toán phân loại, với $p$ là
    tổng số đặc trưng.
\end{enumerate}

Kết quả dự đoán cuối cùng là \emph{majority vote} của tất cả các cây đối với bài
toán phân loại (classification), hoặc \emph{trung bình} đối với hồi quy
(regression).

\begin{figure}[H]
  \centering
  \setlength{\fboxsep}{12pt}
  \fbox{\begin{minipage}{0.88\textwidth}\centering\vspace{1.8cm}
    \FILL{Chèn hình tại đây}\\[0.6em]
    \textit{Gợi ý: Sơ đồ cơ chế Random Forest Ensemble.
    Phía trái: tập dữ liệu gốc. Từ đó vẽ 3--4 mũi tên sang
    các bootstrap sample $D_1, D_2, D_3, D_4$ (mỗi tập có một số bản ghi trùng lặp).
    Mỗi bootstrap sample sinh ra một Decision Tree $h_1, h_2, h_3, h_4$
    (mỗi cây chọn ngẫu nhiên $\sqrt{p}$ đặc trưng tại mỗi nút).
    Phía phải: các dự đoán từng cây được kết hợp qua Majority Vote
    để ra kết quả cuối cùng (Bình thường / Bất thường).}
    \vspace{1.8cm}\end{minipage}}
  \caption{Cơ chế xây dựng và dự đoán của Random Forest Ensemble}
  \label{fig:random-forest}
\end{figure}

\subsubsection{Feature Importance}

Một ưu điểm quan trọng của Random Forest là khả năng đánh giá tầm quan trọng của
từng đặc trưng (feature importance) thông qua chỉ số \emph{Mean Decrease in
Impurity (MDI)}: tổng mức giảm độ không thuần nhất (Gini impurity hoặc entropy) do
đặc trưng đó đóng góp qua tất cả các cây, được chuẩn hóa về tổng bằng 1. Giá trị
\texttt{feature\_importances\_} càng cao, đặc trưng đó càng quan trọng trong việc
phân tách lớp.

\subsubsection{Ứng dụng trong đề tài}

\begin{figure}[H]
  \centering
  \setlength{\fboxsep}{12pt}
  \fbox{\begin{minipage}{0.88\textwidth}\centering\vspace{1.8cm}
    \FILL{Chèn hình tại đây}\\[0.6em]
    \textit{Gợi ý: Biểu đồ cột ngang (horizontal bar chart) thể hiện
    Feature Importance (\texttt{feature\_importances\_}) của mô hình
    Random Forest sau khi huấn luyện trên tập BATADAL.
    Trục X: giá trị importance (0.0--max); Trục Y: tên 31 đặc trưng
    (sắp xếp giảm dần theo importance).
    Tô màu khác nhau cho 3 nhóm đặc trưng: $L_{Ti}$ (xanh dương),
    $F_{PUj}$ (xanh lá), $S_{PUj}$ (cam).
    Hình này được sinh ra từ script \texttt{train\_model.py} sau khi chạy
    \texttt{clf.feature\_importances\_}.}
    \vspace{1.8cm}\end{minipage}}
  \caption{Biểu đồ tầm quan trọng đặc trưng (Feature Importance) của mô hình Random Forest}
  \label{fig:feature-importance}
\end{figure}

Đề tài huấn luyện mô hình \texttt{RandomForestClassifier} (scikit-learn) với
100 cây (\texttt{n\_estimators=100}) và tham số \texttt{class\_weight=`balanced'}
để xử lý mất cân bằng lớp giữa mẫu bình thường (negative) và mẫu tấn công
(positive). Mô hình nhận 31 đặc trưng vật lý ($L_{Ti}$, $F_{PUj}$, $S_{PUj}$,
$S_{V2}$) và dự đoán nhãn nhị phân: 0 (bình thường) hoặc 1 (bất thường). Sau khi
huấn luyện offline, mô hình được serialize thành file \texttt{anomaly\_model.pkl}
và tích hợp vào Spark Structured Streaming thông qua broadcast variable và
\texttt{pandas\_udf} để thực hiện inference theo thời gian thực trên từng
micro-batch.

Đánh giá mô hình sử dụng chiến lược \emph{stratified 75/25 split} và các chỉ số:
Precision, Recall, F1-score và confusion matrix, nhằm đánh giá đúng hiệu năng
trên tập dữ liệu mất cân bằng lớp đặc trưng của bài toán phát hiện bất thường.
